\section{The main component of the pencil configuration space}
\label{sec:main-component}

\Cref{sec:pencil-main-component-route} reduced the hinge-pencil question
to the two main-component statements
(\cref{def:pencil-main-component-statements}): every simple
two-edge-connected multigraph has a pencil realization at the deficiency
rank with adjacent concurrency points distinct, and a generic one
whenever a nondegenerate realization exists at all. This section opens
the geometry those statements are proved from. Over a fixed planar
drawing of the bodies, the pencil realizations of a graph form a linear
space of \emph{heights}; the drawings for which that space is smallest
form a dense open set, over which the realizations assemble into an
irreducible family whose general member attains the deficiency rank. A
single realization at that rank forces the general one, because the
rank is lower semicontinuous. The remaining sections --- the flat rank,
the planar rank input, the ear, split-off, contraction and cut steps,
the structural coverage, and the passage back to the two statements ---
build on the objects introduced here.

\subsection{Planar pictures, the lifting space, and the main component}
\label{sec:main-component-carrier}

Throughout, $G$ is a finite simple connected multigraph of minimum
degree at least two and $K$ is an infinite field.

\begin{definition}[Admissible planar picture]
  \label{def:pencil-admissible-picture}
  \lean{Graph.IsAdmissiblePicture,
        CombinatorialRigidity.Molecular.pencilPicturePoint}
  \leanok
  \uses{def:pencil-nondegenerate}
  A \emph{planar picture} of $G$ is a map $q \colon V(G) \to K^2$,
  written $q_v = (x_v, y_v)$; its \emph{homogeneous picture points} are
  the vectors $(x_v, y_v, 1) \in K^3$. The picture is \emph{admissible}
  if $q_u \neq q_v$ for every edge $uv$ and, for every body $v$, three
  members of the closed neighbourhood of $v$
  (\cref{def:pencil-nondegenerate}) have linearly independent homogeneous
  picture points, that is, the picture points over the closed
  neighbourhood are not collinear.
\end{definition}

Admissibility is a finite conjunction of non-vanishing conditions, one
coordinate difference for each edge and one $3 \times 3$ determinant for
each body, so the admissible pictures form a Zariski-open subset of
$K^{2|V(G)|}$. It is nonempty, since $K$ is infinite and every closed
neighbourhood has at least three members by the degree hypothesis.

\begin{definition}[The lifting space]
  \label{def:pencil-lifting-space}
  \lean{Graph.liftingSpace}
  \leanok
  \uses{def:pencil-admissible-picture}
  Let $q$ be a planar picture of $G$. A \emph{height} is a map
  $z \colon V(G) \to K$. The \emph{lifting space} $L(q)$ is the set of
  heights $z$ such that, for every body $v$, the restriction of $z$ to
  the closed neighbourhood of $v$ agrees with an affine function of the
  picture there: there are $a, b, c \in K$ with
  $z_w = a x_w + b y_w + c$ for every $w$ in that neighbourhood.
\end{definition}

The conditions are linear in $z$, so $L(q)$ is a linear subspace of
$K^{V(G)}$. It contains the space $\mathrm{Aff}(q)$ of globally affine
heights $z_w = \alpha x_w + \beta y_w + \gamma$, which is
three-dimensional when the picture points are not all collinear.

\begin{fmlnote}
  \label{fmlnote:pencil-lifting-space}
  A height is a function on the whole ambient set of bodies that
  vanishes off $V(G)$, so that for a graph not spanning the body set
  the dimension of $L(q)$ is still that of a space of heights on
  $V(G)$.
\end{fmlnote}

\begin{definition}[Configurations and their planes]
  \label{def:pencil-configuration}
  \lean{CombinatorialRigidity.Molecular.pencilConfigPoint,
        CombinatorialRigidity.Molecular.pencilNormalOfPicture}
  \leanok
  \uses{def:pencil-lifting-space}
  A planar picture $q$ and a height $z$ give the \emph{configuration}
  $p = (q, z)$ of homogeneous points
  $p_v = (x_v, y_v, z_v, 1) \in K^4$. Given a choice of three bodies
  $a_v, b_v, c_v$ at every body $v$, the \emph{plane normal} of $p$ at
  $v$ is the vector $n_v \in K^4$ with
  $n_v \cdot w = \det[p_{a_v}, p_{b_v}, p_{c_v}, w]$ for every
  $w \in K^4$.
\end{definition}

The plane normal is polynomial in the picture and height coordinates,
orthogonal to the three chosen points, and nonzero exactly when they
are linearly independent. When the three chosen bodies lie in the
closed neighbourhood of $v$ with non-collinear picture points and
$z \in L(q)$, it is, up to a nonzero scalar, the normal of the plane
containing the points of the whole closed neighbourhood
(\cref{lem:pencil-condition-linear}), whichever three are chosen.

\begin{lemma}[The pencil condition is linear in the heights]
  \label{lem:pencil-condition-linear}
  \uses{def:pencil-lifting-space, def:pencil-configuration,
        def:pencil-panel-realization}
  Let $q$ be admissible and $z$ a height. The configuration $p = (q, z)$
  has every closed neighbourhood coplanar if and only if $z \in L(q)$.
  Every such configuration is a pencil realization
  (\cref{def:pencil-panel-realization}): adjacent bodies carry distinct
  points, and each body's plane is unique and non-vertical.
\end{lemma}
\begin{proof}
  If $z_w = a x_w + b y_w + c$ on the neighbourhood of $v$, its points
  lie on the plane $\zeta = a\xi + b\eta + c$. Conversely a plane through
  those points cannot be vertical --- a vertical plane would cut out an
  affine relation on the non-collinear picture points, which vanish only
  trivially --- so it is the graph of such an affine function, unique
  because the picture points are not collinear. Adjacent points are
  distinct because $q_u \neq q_v$. Linearity is the statement that the
  condition at $v$ places the restriction of $z$ in the image of the
  injective evaluation of affine functions on the neighbourhood, a
  subspace of codimension $\deg v - 2$.
\end{proof}

\begin{lemma}[The hinges are affine in the heights]
  \label{lem:pencil-hinge-affine}
  \uses{def:pencil-configuration, def:pencil-lifting-space, def:extensor}
  Fix a picture $q$. For an edge $e = uv$ the hinge extensor
  $p_u \wedge p_v$ (\cref{def:extensor}) is an affine function of the
  height $z$, its dependence on $z$ linear. Consequently the augmented
  motion matrix at $(q, z)$ splits as a picture-only part plus a part
  linear in $z$; when every hinge extensor is nonzero its rank is
  $|E(G)|$ plus the rank of the panel-hinge rigidity matrix of the
  associated realization. The rank is unchanged when $z$ is scaled by a
  nonzero constant or shifted by a member of $\mathrm{Aff}(q)$.
\end{lemma}
\begin{proof}
  Each Grassmann coordinate of $p_u \wedge p_v$ is a $2\times 2$ minor
  of the two homogeneous points; the coordinates reading the height row
  are linear in $z$, the others constant in $z$. The rank split and the
  scale-and-shift invariance follow by reading the augmented matrix in
  the panel-hinge coordinates, where the height-linear part is supported
  on the hinge columns.
\end{proof}

\begin{definition}[Main pictures]
  \label{def:pencil-main-picture}
  \lean{Graph.IsMainPicture}
  \leanok
  \uses{def:pencil-lifting-space}
  Let $\ell_0$ be the least value of $\dim L(q)$ over admissible
  pictures $q$ of $G$. A \emph{main picture} of $G$ is an admissible
  picture $q$ with $\dim L(q) = \ell_0$, and $U$ denotes the set of main
  pictures.
\end{definition}

\begin{definition}[The general configuration attains]
  \label{def:pencil-x0-attains}
  \lean{Graph.X0Attains}
  \leanok
  \uses{def:pencil-configuration, def:panel-hinge-framework,
        def:rigidity-matrix, def:D-deficiency}
  The \emph{rank} of a configuration $p = (q, z)$ of $G$ is the rank of
  the rigidity matrix (\cref{def:rigidity-matrix}) of the panel-hinge
  framework on $G$ whose panel normal at each body $v$ is $p_v$
  (\cref{def:panel-hinge-framework}). The \emph{general configuration of
  $G$ attains} if there is a nonzero polynomial $P$ in the picture
  coordinates such that every picture $q$ with $P(q) \neq 0$ is
  admissible and the heights $z \in L(q)$ at which $(q, z)$ has rank
  $6(|V(G)| - 1) - \operatorname{def}(\tilde G)$ (\cref{def:D-deficiency},
  at $n = 3$) contain a nonempty Zariski-open subset of $L(q)$.
\end{definition}

The hinge of that framework at an edge $uv$ is the meet of the planes
$p_u^{\perp}$ and $p_v^{\perp}$ (\cref{def:panel-support-extensor}),
the image of the line $p_u \wedge p_v$ under the complement
isomorphism, and it is nonzero whenever $q_u \neq q_v$. An invertible
linear map of the screw space applied to every hinge leaves the rank
unchanged (\cref{lem:screw-map-rows}), so the rank of a configuration
is also the rank of the body-hinge framework with hinges
$p_u \wedge p_v$, the framework of a pencil realization. Taking the
plane normals $n_v$ as panel normals instead would not give this rank:
where two adjacent planes coincide the meet is zero, and a zero hinge
allows no relative motion of its two bodies.

\begin{theorem}[The main component]
  \label{thm:pencil-x0-main-component}
  \uses{def:pencil-main-picture, def:pencil-x0-attains,
        lem:pencil-condition-linear}
  The set $U$ of main pictures (\cref{def:pencil-main-picture}) is
  nonempty and Zariski-open, the pairs $(q, z)$ with $q \in U$ and
  $z \in L(q)$ form a vector bundle over $U$, and its closure $X_0$ is
  irreducible of dimension $2|V(G)| + \ell_0$, containing the flat
  configurations $(q, 0)$ for $q \in U$. The rank of a configuration
  (\cref{def:pencil-x0-attains}) is lower semicontinuous on $X_0$, so
  over the bundle it attains its maximum on a dense open subset. In
  particular, if a single configuration $(q, z)$ with $q \in U$ and
  $z \in L(q)$ has rank at least $6(|V(G)| - 1) - \operatorname{def}(\tilde G)$,
  then the general configuration of $G$ attains, and $G$ has a
  rank-attaining pencil realization with adjacent points distinct.
\end{theorem}
\begin{proof}
  Over an admissible $q$, \cref{lem:pencil-condition-linear} presents
  $L(q)$ as the height-projection of the kernel of the linear system
  fixing the per-body affine functions, a matrix with entries linear in
  $q$; its rank is lower semicontinuous, so the pictures of least
  $\dim L$ are open, and there the kernel is a vector bundle with
  irreducible total space. Over a neighbourhood of a main picture,
  Cramer's rule on a maximal nonvanishing minor gives a polynomial
  section of the bundle through any point of the fibre. Lower
  semicontinuity of the rank on the irreducible $X_0$, together with
  the fact that the rank never exceeds
  $6(|V(G)| - 1) - \operatorname{def}(\tilde G)$ once every hinge is
  nonzero, gives the last two claims; the flat point $(q, 0)$ is a
  member for every $q \in U$.
\end{proof}

\begin{theorem}[The general point of the main component attains]
  \label{thm:pencil-x0-generic-attains}
  \uses{thm:pencil-x0-main-component, def:pencil-x0-attains,
        lem:pencil-hinge-affine, def:pencil-distinct-motive,
        def:pencil-generic-motive, def:D-deficiency,
        def:pencil-main-component-statements}
  For every simple connected multigraph $G$ of minimum degree at least
  two, the general configuration of $G$ attains the rank
  $6(|V(G)| - 1) - \operatorname{def}(\tilde G)$
  (\cref{def:pencil-x0-attains}; \cref{def:D-deficiency}, at $n = 3$).
  Granting this at every such
  graph, the two main-component statements
  (\cref{def:pencil-main-component-statements}) follow: the general
  member of $X_0$ is a pencil realization at the deficiency rank with
  adjacent points distinct (\cref{def:pencil-distinct-motive}), and,
  wherever a nondegenerate realization exists at all, a generic one
  (\cref{def:pencil-generic-motive}).
\end{theorem}
\begin{proof}
  By \cref{thm:pencil-x0-main-component} it suffices to exhibit one
  configuration of the bundle attaining the target rank. That runs a
  strong induction on the number of bodies through ear, split-off,
  contraction and cut steps (the following sections), whose flat inputs
  use the planar rank theorem \cref{thm:theorem-55-6-rows} in place of
  the pin-collinear theorem the informal proof cites. The distinct and
  generic upgrades hold on the dense open subset where adjacent points
  are distinct and, respectively, a nondegenerate realization survives,
  the latter by intersecting two nonempty Zariski-open subsets of one
  fibre $L(q)$;
  the passage on to \cref{thm:pencil-conditional-realization-main-component}
  is the last section.
\end{proof}
